% Arthur Cayley, Collected Mathematical Papers, Volume 7
% Transcription of pages 51-100

\documentclass[12pt,a4paper]{article}

% ==================== PACKAGES ====================
\usepackage[utf8]{inputenc}
\usepackage[T1]{fontenc}
\usepackage[english,french]{babel}
\usepackage{amsmath,amssymb,amsthm}
\usepackage{geometry}
\usepackage{graphicx}
\usepackage{fancyhdr}
\usepackage{titlesec}
\usepackage{enumitem}
\usepackage{array}
\usepackage{booktabs}
\usepackage{mathrsfs}
\usepackage{yfonts}

% ==================== GEOMETRY ====================
\geometry{
  paperwidth=7in,
  paperheight=10in,
  margin=1in,
  top=1in,
  bottom=1in
}

% ==================== PAGE STYLE ====================
\pagestyle{fancy}
\fancyhf{}
\fancyhead[C]{\small\nouppercase\leftmark}
\fancyfoot[C]{\thepage}
\renewcommand{\headrulewidth}{0pt}

% ==================== CUSTOM COMMANDS ====================
\newcommand{\paperref}[1]{\textup{[#1]}}
\newcommand{\article}[1]{\paragraph{#1}}
\DeclareMathOperator{\Det}{Det}

% ==================== THEOREM ENVIRONMENTS ====================
\theoremstyle{definition}
\newtheorem{theorem}{Theorem}
\newtheorem{lemma}{Lemma}
\newtheorem{corollary}{Corollary}

% ==================== DOCUMENT ====================
\begin{document}

% No title page; direct transcription begins
\setcounter{page}{51}

% ============================================================
% PAPER 433: SUR LES SURFACES TETRAEDRALES (continued)
% ============================================================
\selectlanguage{french}

\noindent\paperref{433}\quad
\begin{center}
\textbf{SUR LES SURFACES TETRAEDRALES.}
\end{center}

\noindent
(ce qui donne quatre syst\`emes de valeurs de $\lambda : \mu : \nu$), et puis
\[
\Omega = af\frac{\nu - \mu}{\lambda} + bg\frac{\lambda - \nu}{\mu} + ch\frac{\mu - \lambda}{\nu}:
\]
la surface devient une surface r\'egl\'ee, savoir, la \textit{quadricuspidale} de M.~de~la~Gournerie; et, en supposant
\[
(af)^{\frac{1}{3}} + (bg)^{\frac{1}{3}} + (ch)^{\frac{1}{3}} = 0,
\]
et
\[
\lambda : \mu : \nu = (af)^{\frac{1}{3}} : (bg)^{\frac{1}{3}} : (ch)^{\frac{1}{3}},
\]
ce qui donne
\[
\Omega = \Bigl[(af)^{\frac{1}{3}} - (bg)^{\frac{1}{3}}\Bigr]\,
\Bigl[(bg)^{\frac{1}{3}} - (ch)^{\frac{1}{3}}\Bigr]\,
\Bigl[(ch)^{\frac{1}{3}} - (af)^{\frac{1}{3}}\Bigr]:
\]
la surface devient d\'eveloppable.

\begin{flushright}
\textit{Cambridge, 18 Octobre 1866.}
\end{flushright}

\begin{center}
\textit{Deuxi\`eme M\'emoire. Notes pp.~279---283.}
\end{center}

\smallskip\noindent
\textbf{Note I.}\quad
\textsc{Sur la d\'ecomposition du lieu des g\'en\'eratrices en surfaces t\'etra\'edrales distinctes.}

\smallskip
Il me semble qu'une de vos conclusions a besoin d'\^etre modifi\'ee. Ainsi la surface t\'etra\'edrale d\'eriv\'ee de deux courbes triangulaires \`a exposant $\dfrac{1}{m}$ ($m$ \'etant un entier positif), laquelle, selon un de vos th\'eor\`emes, serait de l'ordre $2m^2$, para\^it se d\'ecomposer en $m$ surfaces chacune de l'ordre $2m$. Il y a pour cela une raison \`a \textit{priori}; en effet, pour deux triangulaires de la forme en question, en employant votre construction, on peut \'etablir une correspondance non-seulement entre les deux syst\`emes de points

% [Illustration: Diagram showing Plan $z=0$ and Plan $w=0$ with curves]

\newpage
\setcounter{page}{52}

\noindent
$A, B, \ldots$, et $A', B', \ldots$, mais aussi entre chaque point $A$ et un seul point correspondant $A'$, car l'\'equation de la premi\`ere courbe \'etant de la forme
\[
fx^{\frac{1}{m}} + gy^{\frac{1}{m}} + hz^{\frac{1}{m}} = 0,
\]
on satisfait \`a cette \'equation en \'ecrivant
\[
x : y : z = a(\theta + \alpha)^m : b(\theta + \beta)^m : c(\theta + \gamma)^m,
\]
o\`u $\theta$ est un param\`etre variable. De m\^eme, l'\'equation de la seconde courbe \'etant
\[
fx^{\frac{1}{m}} + gy^{\frac{1}{m}} + k'w^{\frac{1}{m}} = 0,
\]
on satisfait \`a cette \'equation en \'ecrivant
\[
x : y : w = a'(\theta' + \alpha')^m : b'(\theta' + \beta')^m : d'(\theta' + \delta')^m,
\]
o\`u $\theta'$ est aussi un param\`etre variable.

Pour la droite $OP$, on a
\[
\frac{x}{y} = \frac{a\,(\theta + \alpha)^m}{b\,(\theta + \beta)^m},
\]
et pour la droite $O'P'$
\[
\frac{x}{y} = \frac{a'\,(\theta' + \alpha')^m}{b'\,(\theta' + \beta')^m};
\]
donc, la condition pour la correspondance des droites est
\[
\frac{a\,(\theta + \alpha)^m}{b\,(\theta + \beta)^m}
= \lambda\,\frac{a'\,(\theta' + \alpha')^m}{b'\,(\theta' + \beta')^m},
\]
ce qui donne $m$ valeurs diff\'erentes pour $\theta'$ en termes de $\theta$. Mais chacune de ces valeurs est de la forme
\[
\theta' = \frac{A\theta + B}{C\theta + D},
\]
et, en ne faisant attention qu'\`a une seule valeur de $\theta'$, on a le point
\[
x : y : z = a\,(\theta + \alpha)^m : b\,(\theta + \beta)^m : c\,(\theta + \gamma)^m,
\]
qui correspond \`a un point unique
\[
x : y : w = a'\,(\theta' + \alpha')^m : b'\,(\theta' + \beta')^m : d'\,(\theta' + \delta')^m.
\]

Pour le cas de l'exposant $\frac{1}{3}$, on a, de cette mani\`ere, une surface de l'ordre 6. J'ai v\'erifi\'e cela dans le cas particulier de la surface d\'eveloppable. Il est tr\`es-singulier (c'est M.~Salmon qui m'a fait cette remarque) qu'en \'ecrivant dans cette \'equation $(x^2, y^2, z^2, w^2)$ au lieu de $(x, y, z, w)$, on obtient l'\'equation d'une surface du douzi\`eme ordre, lieu des centres de courbure d'un ellipso"ide.

\begin{flushright}
\textit{Cambridge, 15 F\'evrier 1866.}
\end{flushright}

\newpage
\setcounter{page}{53}

\noindent\paperref{433}\quad
\begin{center}
\textbf{SUR LES SURFACES TETRAEDRALES.}
\end{center}

\smallskip\noindent
\textbf{Note II.}\quad
\textsc{A l'occasion de l'ordre des surfaces t\'etra\'edrales.}

\smallskip
Je crois que j'ai n\'eglig\'e de vous faire conna\^itre un th\'eor\`eme assez g\'en\'eral au sujet de l'ordre de ces surfaces. En consid\'erant dans l'espace deux courbes (planes ou \`a double courbure) des ordres $m$ et $m'$ respectivement, et en supposant qu'il y ait entre les points de ces deux courbes une correspondance $(a, a')$, (c'est-\`a-dire qu'\`a un point donn\'e de la courbe $m$ correspondent $a'$ points sur la courbe $m'$, et \`a un point donn\'e de la courbe $m'$ correspondent $a$ points sur la courbe $m$), alors la surface r\'egl\'ee d\'eriv\'ee est de l'ordre $am' + a'm$. De m\^eme, en consid\'erant dans l'espace deux courbes des ordres $n$ et $n'$ avec une correspondance $(b, b')$, la surface d\'eriv\'ee est de l'ordre $bn' + b'n$; et si l'on a en m\^eme temps les deux syst\`emes de courbes avec les correspondances $(a, a')$ et $(b, b')$, la surface t\'etra\'edrale d\'eriv\'ee est de l'ordre
\[
(am' + a'm)(bn' + b'n) + (an' + a'n)(bm' + b'm).
\]
C'est ce que l'on voit sans peine au moyen des \'equations de la surface t\'etra\'edrale.

En effet, si les \'equations des deux courbes du premier syst\`eme sont
\[
(x, y, z)^m = 0, \quad (x, y, w)^{m'} = 0,
\]
et que la condition de correspondance entre deux points $(x, y, z)$ et $(x, y, w)$ soit exprim\'ee par l'\'equation
\[
(x, y, z, w)^{a, a'} = 0,
\]
o\`u $(x, y, z, w)^{a, a'}$ d\'enote une fonction de l'ordre $a$ en $x, y, z$ et de l'ordre $a'$ en $x, y, w$; alors en \'eliminant $x$ et $y$ entre ces trois \'equations, on obtient pour la surface r\'egl\'ee une \'equation de l'ordre $am' + a'm$ en $z$ et $w$. De m\^eme pour le second syst\`eme de courbes, la surface r\'egl\'ee est de l'ordre $bn' + b'n$. Et pour la surface t\'etra\'edrale, en \'eliminant $x$ et $y$ entre les quatre \'equations, on obtient une \'equation dont l'ordre en $z$ et $w$ est
\[
(am' + a'm)(bn' + b'n) + (an' + a'n)(bm' + b'm),
\]
ce qui est le r\'esultat \'enonc\'e.

\newpage
\setcounter{page}{54}

\noindent
\textbf{Note III.}\quad
\textsc{Sur la surface compl\'ementaire.}

\smallskip
Je profite de cette occasion pour rectifier une erreur que j'ai commise dans ma note sur votre M\'emoire des surfaces t\'etra\'edrales. J'ai dit que la surface compl\'ementaire d'une surface t\'etra\'edrale \'etait une surface r\'egl\'ee; cela est vrai pour certaines surfaces t\'etra\'edrales, mais non pas en g\'en\'eral. Le th\'eor\`eme correct est le suivant:

\textit{La surface compl\'ementaire d'une surface t\'etra\'edrale est la surface t\'etra\'edrale d\'eriv\'ee des deux syst\`emes de courbes, chacun pris deux fois.}

Ainsi, si la surface t\'etra\'edrale est d\'eriv\'ee des courbes $(m, m')$ et $(n, n')$, la surface compl\'ementaire est d\'eriv\'ee des courbes $(m, m)$ et $(n', n')$; et si l'on suppose que la correspondance soit dans chaque cas $(1, 1)$, la surface t\'etra\'edrale sera de l'ordre $2mm' \cdot 2nn' = 4mm'nn'$, et la surface compl\'ementaire sera de l'ordre $2m^2 \cdot 2n'^2 = 4m^2n'^2$; ce qui donne pour la somme des deux surfaces l'ordre $4mm'nn' + 4m^2n'^2$.

\begin{flushright}
\textit{Cambridge, 15 F\'evrier 1866.}
\end{flushright}

\bigskip
\noindent\rule{\textwidth}{0.4pt}

% ============================================================
% PAPER 434: ON CERTAIN SKEW SURFACES, OTHERWISE SCROLLS
% ============================================================
\selectlanguage{english}

\newpage
\setcounter{page}{55}

\noindent\paperref{434}\quad
\begin{center}
\textbf{434.}\\[6pt]
\textbf{ON CERTAIN SKEW SURFACES, OTHERWISE SCROLLS.}\\[4pt]
\small
\textit{[From the \textit{Transactions of the Cambridge Philosophical Society}, vol.~\textsc{xi}. Part~\textsc{ii}. (1869),\\
pp.~277---289. Read Nov.~11, 1867.]}
\end{center}

\normalsize
\setlength{\parindent}{1em}
The investigations contained in the present Memoir were suggested to me by the Memoirs of M.~De~la~Goumerie, presented by him to the Academy of Sciences in 1865 and 1866, published in extract in the Comptes Rendus, and reproduced in his work ``Recherches sur les surfaces reglees tetraedrales symetriques, par Jules De la Goumerie, avec des notes par Arthur Cayley,'' 8vo. Paris, 1867. Although the results or the greater part of them are virtually contained in the formulae established by M.~De~la~Goumerie, yet the special theory of the tetraedral surfaces derived from two pairs of curves of the orders $(m, m')$ and $(n, n')$ respectively, does not appear to have been explicitly developed; and I was led to enter upon it by a very elegant theorem communicated to me by M.~Chasles in September 1865, and subsequently published in the Comptes Rendus, as a portion of the theorem ``Sur les courbes \`a double courbure et sur les surfaces r\'egl\'ees,'' t.~\textsc{lxi}. p.~859 (1865 Oct.~9), viz. if in a tetraedron $ABCD$ we inscribe a variable triangle the sides whereof meet the edges in fixed anharmonic ratios, then the three generating lines of the three hyperboloids $(AB, CD)$, $(AC, BD)$, $(AD, BC)$ respectively, meet in a point; and conversely, if they meet in a point, the anharmonic ratios are equal; or say the variable triangle is inscribed in the tetraedron with a given form; the locus of the point of intersection is a skew surface or \textit{scroll}, the order of which was to be determined.

The particular case of Chasles's theorem, where the variable triangle reduces itself to a point, gives the theorem that the locus of a point such that its projections from the four edges on the opposite faces are generating lines of the hyperboloids, is a sextic surface, the locus of the point of intersection of the three generating lines: this surface was mentioned as a sextic scroll by M.~Chasles, and the same surface was considered by me in the paper ``On Tetrahedral Surfaces,'' Phil.~Trans. vol.~\textsc{cliii}. (1863), pp.~417--436, see p.~420.

\newpage
\setcounter{page}{56}

\noindent\paperref{434}\quad
\textbf{ON CERTAIN SKEW SURFACES, OTHERWISE SCROLLS.}

\smallskip
The Memoir is divided into three Parts.

\textsc{Part I.}\quad Article Nos. 1 to 9. On the scrolls in general.

\textsc{Part II.}\quad Article Nos. 10 to 14. On the particular class of scrolls.

\textsc{Part III.}\quad Article Nos. 15 to 29. The Curves $U$, $U'$ are henceforward ``triangular'' curves.

\smallskip
\textbf{Article Nos.~1 to 9. On the scrolls in general.}

\smallskip
\textit{Generation of the scroll.}

1.~Consider any two curves $U$, $U'$ (in space, planes or skew); let their orders be $m$, $m'$ respectively, and let there be between the points of the two curves a correspondence $(\alpha, \alpha')$, viz. to each point of $U$ correspond $\alpha'$ points of $U'$, and to each point of $U'$ correspond $\alpha$ points of $U$; the number of united points (or points corresponding to each other) on $U$ is $= \alpha + \alpha'$; and in like manner the number on $U'$ is $= \alpha + \alpha'$.

2.~In particular let $U$, $U'$ be the sections of a scroll $S$ by two fixed planes; the scroll is generated by the lines which join the corresponding points of the two curves; the class of generation is $(\alpha, \alpha')$; and the order of the scroll is $= m\alpha' + m'\alpha$. This includes as a particular case the generation of a scroll by means of two director curves and a director line: viz. if the director line meet the planes of $U$, $U'$ in the points $O$, $O'$ respectively, and if $OP$, $O'P'$ are corresponding generating lines, then the points $P$, $P'$ on the director curves have a correspondence $(1, 1)$, and the order of the scroll is $= m + m'$.

3.~This conclusion may be otherwise established by means of the section of the scroll by an arbitrary plane: viz. taking the section by the plane of $U$, this is made up of the curve $U$ taken $\alpha'$ times (or, what is the same thing, the complete section is a curve of the order $m\alpha'$), together with the $\alpha$ lines through the several intersections of the plane with the director line; the total order of the section is $= m\alpha' + \alpha$; but this is equal to the order of the scroll into the number of generating lines through an arbitrary point of the plane of $U$: an arbitrary point of the plane has with the curve $U$ a $(1, m)$ correspondence, and therefore with the curve $U'$ a $(\alpha, m\alpha')$ correspondence; hence through the point pass $m\alpha'$ generating lines of the scroll, and the order of the scroll is thus $= m\alpha' + m'\alpha$.

\newpage
\setcounter{page}{57}

\noindent\paperref{434}\quad
\textbf{ON CERTAIN SKEW SURFACES, OTHERWISE SCROLLS.}

4.~There are on the scroll certain curves $V$ of the order $m$, viz. the section by the plane of $U$ consists of the curve $U$ ($= m$) together with the generating lines through the $m$ intersections of the plane with the director line; each of these lines meets $U$ in a single point, and the $m$ points thus obtained are points of a curve $V$ of the order $m$. There are in like manner curves $V'$ of the order $m'$.

5.~The number of these curves $V$ is $= \alpha$; and the like number of curves $V'$ is $= \alpha'$.

6.~In particular the curves $U$, $U'$ may reduce themselves each to a point: we have then a scroll of the order 2, viz. this is a quadric surface, which may be considered as a scroll generated by a variable line which meets each of two given lines, and also each of two given points; or say the scroll has for its director lines two lines, and for its director points two points.

\textit{Section of the scroll by the plane of $U$ or of $U'$.}

7.~By what precedes, the scroll $S$ ($m\alpha' + m'\alpha$) meets the plane of $U$ in the curve $U$ taken $\alpha'$ times (order $m\alpha'$) together with $\alpha m'$ generating lines; and it meets the plane of $U'$ in the curve $U'$ taken $\alpha$ times (order $m'\alpha$) together with $\alpha' m$ generating lines.

\textit{Section of the scroll by the tangent plane at a point of $U$ or of $U'$.}

8.~The foregoing results will subsist if for the arbitrary plane we substitute the tangent plane at a point $A$ of the curve $U$: the section by the tangent plane at $A$ is made up of the curve $U$ ($= m$) taken $\alpha'$ times, order $m\alpha'$, together with the generating lines through the $m'$ intersections (other than $A$) of the tangent plane with the director line. But each of these generating lines meets $U$ in $\alpha$ points (besides the point $A$), and there are thus $\alpha m' - \alpha$ of these points, giving rise to $\alpha m' - \alpha$ generating lines which lie in the tangent plane; the complete section by the tangent plane is therefore of the order $m\alpha' + \alpha m' - \alpha$, that is, $= (m + m')\alpha' - \alpha$, or say $= m\alpha' + m'\alpha - \alpha$, which is the order of the scroll, and the section is therefore the complete section.

\newpage
\setcounter{page}{58}

\noindent\paperref{434}\quad
\textbf{ON CERTAIN SKEW SURFACES, OTHERWISE SCROLLS.}

9.~The section by the tangent plane through $A$ is made up of $(m - \alpha)(m' - \alpha') - \alpha\alpha'$ generating lines (viz. these are, the line $V'\!A$ $\alpha'(m - \omega - \alpha)$ times, the line $V\!A$ $\alpha(m' - \omega' - \alpha')$ times, and the line $AA'$ $\alpha\alpha'$ times, if $\omega$, $\omega'$ are certain integers), and of the curve $U$ taken $\alpha'$ times (order $m\alpha'$), and the curve $U'$ taken $\alpha$ times (order $m'\alpha$).

\bigskip
\textbf{Article Nos.~10 to 14. On the particular class of scrolls.}

\smallskip
\textit{10. Geometrical Construction of a class of Scrolls.}

1.~Consider any two curves $\Omega$, $\Omega'$ of the same deficiency, and such that to each point of $\Omega$ there corresponds a single tangent of $\Omega'$, and to each tangent of $\Omega'$ a single tangent of $\Omega$ (this assumes that the curves $\Omega$, $\Omega'$ are rational transformations one of the other, and that they have consequently the same Deficiency). This being so, let the points of $U$ which lie on any tangent of $\Omega$ and the points of $U'$ which lie on the corresponding tangent of $\Omega'$ be taken to be corresponding points of $U$, $U'$.

The correspondence is then $(\mu\mu', \mu'\mu)$; in fact through a given point of $U$ there pass $\mu$ tangents of $\Omega$, and the corresponding $\mu'$ tangents of $\Omega'$ each intersect $U'$ in a single point $P'$; that is, to a given point $P$ correspond $\mu'$ points $P'$; and similarly to a given point $P'$ correspond $\mu$ points $P$.

And this being so, the number of united points, that is, points of $U$ through which pass corresponding tangents of $\Omega$, $\Omega'$, is $= \mu + \mu'$.

\newpage
\setcounter{page}{59}

\noindent\paperref{434}\quad
\textbf{ON CERTAIN SKEW SURFACES, OTHERWISE SCROLLS.}

2.~The generation of the scroll is obvious: to each point $P$ of $U$ we have the corresponding tangent of $\Omega$, and taking on this the $\alpha'$ corresponding points of $U'$, we have through $P$ the $\alpha'$ generating lines of the scroll; and the order of the scroll is $= m\alpha' + m'\alpha$.

3.~This conclusion may be otherwise established as follows: the scroll is generated by the lines which join corresponding points of $U$ and $U'$; these lines meet each of the planes of $U$, $U'$ in a single point, and therefore meet any plane whatever in a single point; that is, they are generating lines of a scroll. And the order of the scroll is $= m\alpha' + m'\alpha$, as appears by the section made by the plane of $U$.

4.~There are on the scroll curves $V$ of the order $m$, and curves $V'$ of the order $m'$; the number of curves $V$ is $= \alpha$, and the number of curves $V'$ is $= \alpha'$.

5.~The section by the plane of $U$ is made up of the curve $U$ taken $\alpha'$ times (order $m\alpha'$), together with $\alpha m'$ generating lines.

6.~In particular the curves $\Omega$, $\Omega'$ may reduce themselves each to a point: the scroll is then a quadric surface, which may be considered as generated by a variable line meeting each of two given lines and each of two given points.

7.~By what precedes, the scroll $S$ ($m\alpha' + m'\alpha$) meets the plane of $U$ in the curve $U$ taken $\alpha'$ times (order $m\alpha'$) together with $\alpha m'$ generating lines; and it meets the plane of $U'$ in the curve $U'$ taken $\alpha$ times (order $m'\alpha$) together with $\alpha' m$ generating lines.

8.~The section by the tangent plane through $A$ is made up of the curve $U$ taken $\alpha'$ times (order $m\alpha'$), together with $\alpha m' - \alpha$ generating lines, and is of the order $m\alpha' + \alpha m' - \alpha = m\alpha' + m'\alpha - \alpha$, which is the order of the scroll.

\newpage
\setcounter{page}{60}

\noindent\paperref{434}\quad
\textbf{ON CERTAIN SKEW SURFACES, OTHERWISE SCROLLS.}

9.~The section by the tangent plane through $A$ is made up of $(m - \alpha)(m' - \alpha') - \alpha\alpha'$ generating lines (viz. these are, the line $V'\!A$ $\alpha'(m - \omega - \alpha)$ times, the line $V\!A$ $\alpha(m' - \omega' - \alpha')$ times, and the line $AA'$ $\alpha\alpha'$ times), and of the curve $U$ taken $\alpha'$ times (order $m\alpha'$), and the curve $U'$ taken $\alpha$ times (order $m'\alpha$).

\bigskip
\textbf{Article Nos.~15 to 29. The Curves $U$, $U'$ are henceforward ``triangular'' curves.}

\smallskip
\textit{15.}
Let $r = \pm\dfrac{p}{q}$, where $p$, $q$ are positive integers prime to each other, and let the given sections be
\begin{align*}
z &= 0, \quad Ax^r + By^r \quad\; + Dw^r = 0,\\
w &= 0, \quad A'x^r + B'y^r + C'z^r \quad\; = 0,
\end{align*}
where it is to be observed that $r$ being $= +\dfrac{p}{q}$, the two given sections are of the order $pq$, the order of the scroll is $= 2p^2q^2$; each of the given sections is a $pq$-tuple line on the scroll, and the plane thereof meets the scroll in the section taken $pq$ times, and in the $pq$ generating lines: but $r$ being $= -\dfrac{p}{q}$, the two given sections are each of the order $2pq$, with three $pq$-tuple points ($\omega = \alpha = \beta = pq$, $\omega' = \alpha' = \beta' = pq$), and thence the order of the scroll is $\frac{1}{2}(2pq)^2 = 2p^2q^2$; each of the sections is a $pq$-tuple line on the scroll, and the plane meets the scroll only in the section taken $pq$ times. But in either case, if $q$ be $> 1$, that is, if $r$ be fractional, it will presently appear that the scroll of the order $2p^2q^2$ breaks up into $q$ scrolls each of the order $2p^2q$.

\newpage
\setcounter{page}{61}

\noindent\paperref{434}\quad
\textbf{ON CERTAIN SKEW SURFACES, OTHERWISE SCROLLS.}

\textit{16. Analytical Theory of the scrolls.}

11.~It will be convenient to consider first the case $r = +\dfrac{p}{q}$; and to simplify the formulae I write
\[
\xi = x^{\frac{1}{q}}, \quad \eta = y^{\frac{1}{q}}, \quad \zeta = z^{\frac{1}{q}}, \quad \omega = w^{\frac{1}{q}}.
\]
The equations of the two curves $U$, $U'$ then are
\begin{align*}
\zeta^q &= 0, \quad A\xi^p + B\eta^p + D\omega^p = 0,\\
\omega^q &= 0, \quad A'\xi^p + B'\eta^p + C'\zeta^p = 0;
\end{align*}
and the equations of any two corresponding points thereof may be taken to be
\begin{alignat*}{2}
(x, y, z, w) &= (\alpha(\theta + \alpha)^q,\; &&\beta(\theta + \beta)^q,\; \gamma(\theta + \gamma)^q,\; 0),\\
(x, y, z, w) &= (\alpha'(\theta' + \alpha')^q,\; &&\beta'(\theta' + \beta')^q,\; 0,\; \delta'(\theta' + \delta')^q),
\end{alignat*}
respectively; whence the equations of the line joining the two points are
\begin{align*}
\frac{x}{\alpha(\theta + \alpha)^q} &= \frac{y}{\beta(\theta + \beta)^q} = \frac{z}{\gamma(\theta + \gamma)^q},\\
\frac{x}{\alpha'(\theta' + \alpha')^q} &= \frac{y}{\beta'(\theta' + \beta')^q} = \frac{w}{\delta'(\theta' + \delta')^q};
\end{align*}
that is,
\[
\frac{x}{\alpha(\theta + \alpha)^q} = \frac{y}{\beta(\theta + \beta)^q},
\qquad
\frac{z}{\gamma(\theta + \gamma)^q} = 0,
\]
and
\[
\frac{x}{\alpha'(\theta' + \alpha')^q} = \frac{y}{\beta'(\theta' + \beta')^q},
\qquad
\frac{w}{\delta'(\theta' + \delta')^q} = 0.
\]

\newpage
\setcounter{page}{62}

\noindent\paperref{434}\quad
\textbf{ON CERTAIN SKEW SURFACES, OTHERWISE SCROLLS.}

12.~Or, what is the same thing, the line is the intersection of the two planes
\begin{align*}
\frac{x}{\alpha(\theta + \alpha)^q} - \frac{y}{\beta(\theta + \beta)^q} &= 0,
& \frac{z}{\gamma(\theta + \gamma)^q} &= 0,\\
\frac{x}{\alpha'(\theta' + \alpha')^q} - \frac{y}{\beta'(\theta' + \beta')^q} &= 0,
& \frac{w}{\delta'(\theta' + \delta')^q} &= 0.
\end{align*}

13.~Substituting the values $X = \theta Y$, $X' = \theta' Y'$, we have
\[
M(\theta Y, Y, W) = 0, \quad M(\theta' Y', Y', W') = 0,
\]
and thence eliminating $\theta$, $\theta'$, we obtain the equation of the scroll.

14.~We may find the sections of the scroll by the coordinate planes, and the sections by the planes through the coordinate axes, without obtaining the equation of the scroll itself. Thus the section by the plane $z = 0$ is found by eliminating $\theta$, $\theta'$ between the three equations
\[
\frac{x}{\alpha(\theta + \alpha)^q} = \frac{y}{\beta(\theta + \beta)^q},
\qquad
\frac{x}{\alpha'(\theta' + \alpha')^q} = \frac{y}{\beta'(\theta' + \beta')^q},
\qquad
\frac{w}{\delta'(\theta' + \delta')^q} = 0.
\]

\newpage
\setcounter{page}{63}

\noindent\paperref{434}\quad
\textbf{ON CERTAIN SKEW SURFACES, OTHERWISE SCROLLS.}

15.~The section by the plane $z = 0$ is thus a curve of the order $2pq$, viz. it is composed of the two curves
\begin{align*}
w &= 0, \quad Ax^r + By^r = 0,\\
w &= 0, \quad A'x^r + B'y^r = 0,
\end{align*}
each taken $q$ times.

16.~In like manner the section by the plane $w = 0$ is a curve of the order $2pq$, composed of the two curves
\begin{align*}
z &= 0, \quad Ax^r + By^r = 0,\\
z &= 0, \quad A'x^r + B'y^r = 0,
\end{align*}
each taken $q$ times.

17.~The section by the plane $x = 0$ is found by eliminating $\theta$, $\theta'$ between the three equations
\[
\frac{1}{\alpha(\theta + \alpha)^q} = \frac{y}{\beta(\theta + \beta)^q},
\qquad
\frac{1}{\alpha'(\theta' + \alpha')^q} = \frac{y}{\beta'(\theta' + \beta')^q},
\qquad
\frac{z}{\gamma(\theta + \gamma)^q} = 0,\;
\frac{w}{\delta'(\theta' + \delta')^q} = 0.
\]

18.~The section by the plane $x = 0$ is a curve of the order $2pq$, composed of the curves
\begin{align*}
y &= 0, \quad C'z^r + Dw^r = 0,\\
y &= 0, \quad Cz^r + D'w^r = 0,
\end{align*}
each taken $q$ times.

\newpage
\setcounter{page}{64}

\noindent\paperref{434}\quad
\textbf{ON CERTAIN SKEW SURFACES, OTHERWISE SCROLLS.}

19.~The section by the plane $y = 0$ is a curve of the order $2pq$, composed of the curves
\begin{align*}
x &= 0, \quad C'z^r + Dw^r = 0,\\
x &= 0, \quad Cz^r + D'w^r = 0,
\end{align*}
each taken $q$ times.

20.~The section by the plane through the axis $z$, $w$ (that is, the plane $x = 0$, $y = 0$) is a curve of the order $2pq$, composed of the two lines $z = 0$, $w = 0$ each taken $pq$ times.

21.~The section by the plane through the axis $x$, $y$ (that is, the plane $z = 0$, $w = 0$) is a curve of the order $2pq$, composed of the two lines $x = 0$, $y = 0$ each taken $pq$ times.

22.~The section by the plane through the axis $x$, $z$ (that is, the plane $y = 0$, $w = 0$) is a curve of the order $2pq$, composed of the two lines $x = 0$, $z = 0$ each taken $pq$ times.

23.~The section by the plane through the axis $y$, $w$ (that is, the plane $x = 0$, $z = 0$) is a curve of the order $2pq$, composed of the two lines $y = 0$, $w = 0$ each taken $pq$ times.

24.~It may be remarked that the sections by the planes $z = 0$ and $w = 0$ each consist of two curves of the order $pq$, each taken $q$ times; and that the sections by the planes $x = 0$ and $y = 0$ each consist of two curves of the order $pq$, each taken $q$ times.

25.~The sections by the planes through the coordinate axes each consist of two lines, each taken $pq$ times.

26.~We see that the scroll is a surface of the order $2p^2q^2$, having the curves $U$ and $U'$ each as $pq$-tuple lines; and the planes of these curves meet the scroll in these curves taken $pq$ times, and in $pq$ generating lines.

\newpage
\setcounter{page}{65}

\noindent\paperref{434}\quad
\textbf{ON CERTAIN SKEW SURFACES, OTHERWISE SCROLLS.}

27.~If $q > 1$, that is, if $r$ be fractional, the scroll breaks up into $q$ scrolls, each of the order $2p^2q$; and each of these scrolls has the curves $U$ and $U'$ each as $p$-tuple lines, and the planes of these curves meet each scroll in these curves taken $p$ times, and in $p$ generating lines.

28.~The case $r = -\dfrac{p}{q}$ is in like manner obtained by writing in the formulae for $r = +\dfrac{p}{q}$, $x^{-1}$, $y^{-1}$, $z^{-1}$, $w^{-1}$ in place of $x$, $y$, $z$, $w$ respectively. The resulting scroll is of the same order $2p^2q^2$, but the nature of the singularities is different.

29.~The scrolls of the order $2p^2q$ are generated by lines which join corresponding points of two curves $U$, $U'$; and the correspondence is established by means of two curves $\Omega$, $\Omega'$ of the same deficiency, such that to each point of $\Omega$ corresponds a single tangent of $\Omega'$, and to each tangent of $\Omega'$ corresponds a single tangent of $\Omega$.

\bigskip
\begin{flushright}
\textit{Cambridge, Nov.~11, 1867.}
\end{flushright}

\bigskip
\noindent\rule{\textwidth}{0.4pt}

% ============================================================
% PAPER 435: ON THE SIX COORDINATES OF A LINE
% ============================================================
\newpage
\setcounter{page}{66}

\noindent\paperref{435}\quad
\begin{center}
\textbf{435.}\\[6pt]
\textbf{ON THE SIX COORDINATES OF A LINE.}\\[4pt]
\small
\textit{[From the \textit{Transactions of the Cambridge Philosophical Society}, vol.~\textsc{xi}. Part~\textsc{ii}. (1871),\\
pp.~231---238. Read March~9, 1868.]}
\end{center}

\normalsize
\setlength{\parindent}{1em}
\textit{Introduction. Article Nos.~1 to 8.}

1.~In the geometry of three dimensions, a line may be determined by means of six coordinates $(a, b, c, f, g, h)$, which are not independent, but are connected by the equation $af + bg + ch = 0$. These coordinates may be regarded as the coefficients of a bilinear form, or as the elements of an anti-symmetric matrix; and they serve to express the condition that two lines may intersect, or that a line may belong to a given complex, congruence, or regulus. The coordinates were first introduced by Grassmann in 1844, and independently by Pl"ucker in 1846; and they have been extensively employed by Cayley, Salmon, and other writers.

2.~If $(x, y, z, w)$ are the coordinates of a point, and $(X, Y, Z, W)$ the coordinates of a plane, then the line joining the two points $(x, y, z, w)$ and $(x', y', z', w')$ has for its six coordinates
\begin{align*}
a &= y z' - y' z, & b &= z x' - z' x, & c &= x y' - x' y,\\
f &= x w' - x' w, & g &= y w' - y' w, & h &= z w' - z' w;
\end{align*}
and these satisfy the relation $af + bg + ch = 0$.

\newpage
\setcounter{page}{67}

\noindent\paperref{435}\quad
\textbf{ON THE SIX COORDINATES OF A LINE.}

3.~The six coordinates may also be expressed in terms of the coordinates of two planes. If $(X, Y, Z, W)$ and $(X', Y', Z', W')$ are the coordinates of two planes, then their line of intersection has for its six coordinates
\begin{align*}
a &= X W' - X' W, & b &= Y W' - Y' W, & c &= Z W' - Z' W,\\
f &= Y Z' - Y' Z, & g &= Z X' - Z' X, & h &= X Y' - X' Y;
\end{align*}
and these also satisfy the relation $af + bg + ch = 0$.

4.~The condition that two lines $(a, b, c, f, g, h)$ and $(a_1, b_1, c_1, f_1, g_1, h_1)$ may intersect is
\[
a f_1 + b g_1 + c h_1 + f a_1 + g b_1 + h c_1 = 0.
\]
This may be written in the form
\[
\begin{vmatrix}
a & b & c & f & g & h \\
a_1 & b_1 & c_1 & f_1 & g_1 & h_1
\end{vmatrix}
= 0,
\]
where the expression on the left denotes the sum of the products of corresponding elements, taken with the proper signs.

5.~If the two lines coincide, we have the conditions
\[
\frac{a}{a_1} = \frac{b}{b_1} = \frac{c}{c_1} = \frac{f}{f_1} = \frac{g}{g_1} = \frac{h}{h_1},
\]
or, what is the same thing,
\[
a : b : c : f : g : h = a_1 : b_1 : c_1 : f_1 : g_1 : h_1.
\]

\newpage
\setcounter{page}{68}

\noindent\paperref{435}\quad
\textbf{ON THE SIX COORDINATES OF A LINE.}

6.~The condition that the line $(a, b, c, f, g, h)$ may belong to the linear complex
\[
AX + BY + CZ + DW + EU + FV = 0
\]
is
\[
Aa + Bb + Cc + Df + Eg + Fh = 0.
\]
The coordinates of the line satisfy this linear equation; and conversely, any six quantities satisfying this equation may be taken as the coordinates of a line belonging to the complex.

7.~The condition that the line $(a, b, c, f, g, h)$ may belong to the linear congruence
\begin{align*}
A_1 X + B_1 Y + C_1 Z + D_1 W + E_1 U + F_1 V &= 0,\\
A_2 X + B_2 Y + C_2 Z + D_2 W + E_2 U + F_2 V &= 0,
\end{align*}
is that its coordinates satisfy both of the linear equations
\begin{align*}
A_1 a + B_1 b + C_1 c + D_1 f + E_1 g + F_1 h &= 0,\\
A_2 a + B_2 b + C_2 c + D_2 f + E_2 g + F_2 h &= 0.
\end{align*}

8.~The condition that the line $(a, b, c, f, g, h)$ may belong to the regulus
\[
\Phi(X, Y, Z, W, U, V) = 0
\]
is that its coordinates satisfy the quadratic equation
\[
\Phi(a, b, c, f, g, h) = 0.
\]

\newpage
\setcounter{page}{69}

\noindent\paperref{435}\quad
\textbf{ON THE SIX COORDINATES OF A LINE.}

formulae of transformation to a new set of coordinate planes $x_0 = 0$, $y_0 = 0$, $z_0 = 0$, $w_0 = 0$. And I shall also show in what manner the absolute magnitudes of the coordinates may be fixed. But deferring the consideration of these questions, I consider the planes $x = 0$, $y = 0$, $z = 0$, $w = 0$ as given planes, and take the six coordinates $(a, b, c, f, g, h)$ of a line to be determined as above in reference to these given planes, the absolute values of these coordinates remaining indeterminate, and their ratios only being attended to. And I proceed to consider the various questions which present themselves in the geometry of the line, considered as thus determined by means of its six coordinates $(a, b, c, f, g, h)$.

\smallskip
\textbf{Article, Nos.~9 to 18. (Various Sub-headings.) Elementary Theorems.}

\smallskip
\textit{Condition that a line may be in a given plane.}

9.~Taking the line to be $(a, b, c, f, g, h)$, the equation of the given plane to be
\[
(A, B, C, D \,\backslash\!\!\backslash\, x, y, z, w) = 0;
\]
then if $(\alpha, \beta, \gamma, \delta)$, $(\alpha', \beta', \gamma', \delta')$ are the coordinates of any two points on the line, we have the system of equations \textit{ante}, No.~2, and substituting therein for $\beta\gamma' - \beta'\gamma$, \&c. the values $(a, b, c, f, g, h)$, we find
\[
\begin{vmatrix}
0, & c, & -b, & f \\
-c, & 0, & a, & g \\
b, & -a, & 0, & h \\
-f, & -g, & -h, & 0
\end{vmatrix}
(A, B, C, D) = 0;
\]
which equations, equivalent to a twofold relation, are the required condition. It may be remarked that, treating $(A, B, C, D)$ as current plane coordinates, each equation of the system is that of a point lying in the line.

\newpage
\setcounter{page}{70}

\noindent\paperref{435}\quad
\textbf{ON THE SIX COORDINATES OF A LINE.}

\textit{Condition that a line may pass through a given point.}

10.~The coordinates of the given point are taken to be $(\alpha, \beta, \gamma, \delta)$. If
\[
(A, B, C, D \,\backslash\!\!\backslash\, x, y, z, w) = 0, \quad
(A', B', C', D' \,\backslash\!\!\backslash\, x, y, z, w) = 0,
\]
are the equations of any two planes through the line, then we have the system of equations \textit{ante} No.~3, and substituting therein for $AB' - A'B$, \&c. their values in terms of the coordinates $(a, b, c, f, g, h)$ of the line, we have
\[
\begin{vmatrix}
0, & h, & -g, & a \\
-h, & 0, & f, & b \\
g, & -f, & 0, & c \\
-a, & -b, & -c, & 0
\end{vmatrix}
(\alpha, \beta, \gamma, \delta) = 0;
\]
which equations, equivalent to a twofold relation, are the required condition. It is obvious that, treating $(\alpha, \beta, \gamma, \delta)$ as current point coordinates, each equation of the system is the equation of a plane through the given line.

\textit{Condition that two lines may intersect.}

11.~For two lines $(a, b, c, f, g, h)$ and $(a_1, b_1, c_1, f_1, g_1, h_1)$ the condition of intersection is, as already mentioned,
\[
a f_1 + b g_1 + c h_1 + f a_1 + g b_1 + h c_1 = 0.
\]

\newpage
\setcounter{page}{71}

\noindent\paperref{435}\quad
\textbf{ON THE SIX COORDINATES OF A LINE.}

It may be remarked that the expression on the left-hand side is the determinant of the matrix
\[
\begin{pmatrix}
a & b & c & f & g & h \\
a_1 & b_1 & c_1 & f_1 & g_1 & h_1
\end{pmatrix},
\]
and that the vanishing of this determinant expresses that the two lines intersect. The determinant may be written in the form
\[
\begin{vmatrix}
a & h & -g & a_1 & h_1 & -g_1 \\
-h & b & f & -h_1 & b_1 & f_1 \\
g & -f & c & g_1 & -f_1 & c_1 \\
a & b & c & a_1 & b_1 & c_1
\end{vmatrix}
= 0,
\]
or in other equivalent forms.

\textit{Condition that a line may meet each of four given lines.}

12.~If the four given lines are $(a_1, b_1, c_1, f_1, g_1, h_1)$, $(a_2, b_2, c_2, f_2, g_2, h_2)$, $(a_3, b_3, c_3, f_3, g_3, h_3)$, $(a_4, b_4, c_4, f_4, g_4, h_4)$, then the conditions that the line $(a, b, c, f, g, h)$ may meet each of them are
\begin{align*}
a f_1 + b g_1 + c h_1 + f a_1 + g b_1 + h c_1 &= 0,\\
a f_2 + b g_2 + c h_2 + f a_2 + g b_2 + h c_2 &= 0,\\
a f_3 + b g_3 + c h_3 + f a_3 + g b_3 + h c_3 &= 0,\\
a f_4 + b g_4 + c h_4 + f a_4 + g b_4 + h c_4 &= 0.
\end{align*}
These are four linear equations in the six coordinates $(a, b, c, f, g, h)$, which together with the quadratic relation $af + bg + ch = 0$ determine the ratios of these coordinates. The number of solutions is in general $= \infty^1$, that is, the lines meeting four given lines form a regulus.

\newpage
\setcounter{page}{72}

\noindent\paperref{435}\quad
\textbf{ON THE SIX COORDINATES OF A LINE.}

\textit{Condition that a line may belong to a given complex of the second order.}

13.~The general equation of a complex of the second order is
\[
(A, B, C, F, G, H \,\backslash\!\!\backslash\, a, b, c, f, g, h)^2 = 0,
\]
where $(A, B, C, F, G, H \,\backslash\!\!\backslash\, a, b, c, f, g, h)^2$ denotes a quadratic function of the six coordinates. The condition that the line $(a, b, c, f, g, h)$ may belong to the complex is that its coordinates satisfy this equation.

14.~The complex of the second order may be represented as the assemblage of lines which satisfy the equation
\[
\lambda(a, b, c, f, g, h)^2 + \mu(a, b, c, f, g, h)(a_1, b_1, c_1, f_1, g_1, h_1) + \nu(a_1, b_1, c_1, f_1, g_1, h_1)^2 = 0,
\]
where $\lambda$, $\mu$, $\nu$ are parameters, and $(a, b, c, f, g, h)$, $(a_1, b_1, c_1, f_1, g_1, h_1)$ are the coordinates of two fixed lines.

\textit{Condition that a line may belong to a given congruence of the second order.}

15.~The general equation of a congruence of the second order is given by the intersection of two complexes of the second order; the coordinates of the lines belonging to the congruence satisfy two quadratic equations.

\newpage
\setcounter{page}{73}

\noindent\paperref{435}\quad
\textbf{ON THE SIX COORDINATES OF A LINE.}

16.~The congruence of the second order may be represented as the assemblage of lines which satisfy the equations
\begin{align*}
&(A_1, B_1, C_1, F_1, G_1, H_1 \,\backslash\!\!\backslash\, a, b, c, f, g, h)^2 = 0,\\
&(A_2, B_2, C_2, F_2, G_2, H_2 \,\backslash\!\!\backslash\, a, b, c, f, g, h)^2 = 0,
\end{align*}
where the two quadratic forms are independent.

\textit{Condition that a line may belong to a given regulus.}

17.~The regulus is a special case of the congruence; it is the assemblage of lines which meet each of three given lines. The condition that the line $(a, b, c, f, g, h)$ may belong to a regulus is that its coordinates satisfy three linear equations, which express that the line meets each of the three given lines.

18.~Conversely, any three linear equations in the six coordinates determine a regulus; for they express that the line meets each of three given lines, and the lines meeting three given lines form a regulus.

\newpage
\setcounter{page}{74}

\noindent\paperref{435}\quad
\textbf{ON THE SIX COORDINATES OF A LINE.}

\smallskip
\textbf{Article, Nos.~19 to 35. Various Formulae.}

\smallskip
\textit{Expression for the coordinates of a line in terms of the coordinates of two points.}

19.~If $(x, y, z, w)$ and $(x', y', z', w')$ are two points on the line, then the six coordinates of the line are given by the formulae
\begin{align*}
a &= y z' - y' z, & b &= z x' - z' x, & c &= x y' - x' y,\\
f &= x w' - x' w, & g &= y w' - y' w, & h &= z w' - z' w;
\end{align*}
and conversely, if $(a, b, c, f, g, h)$ are the coordinates of a line, and $(x, y, z, w)$ is a point on the line, then the coordinates of any other point on the line may be expressed in terms of these quantities.

20.~If $(X, Y, Z, W)$ and $(X', Y', Z', W')$ are the coordinates of two planes through the line, then the six coordinates of the line are given by the formulae
\begin{align*}
a &= X W' - X' W, & b &= Y W' - Y' W, & c &= Z W' - Z' W,\\
f &= Y Z' - Y' Z, & g &= Z X' - Z' X, & h &= X Y' - X' Y.
\end{align*}

\newpage
\setcounter{page}{75}

\noindent\paperref{435}\quad
\textbf{ON THE SIX COORDINATES OF A LINE.}

\textit{Expression for the coordinates of a line in terms of the coordinates of two planes.}

21.~If $(X, Y, Z, W)$ and $(X', Y', Z', W')$ are the coordinates of two planes through the line, then the six coordinates of the line may be expressed in the form
\begin{align*}
a &= \begin{vmatrix} X & W \\ X' & W' \end{vmatrix},
& b &= \begin{vmatrix} Y & W \\ Y' & W' \end{vmatrix},
& c &= \begin{vmatrix} Z & W \\ Z' & W' \end{vmatrix},\\
f &= \begin{vmatrix} Y & Z \\ Y' & Z' \end{vmatrix},
& g &= \begin{vmatrix} Z & X \\ Z' & X' \end{vmatrix},
& h &= \begin{vmatrix} X & Y \\ X' & Y' \end{vmatrix}.
\end{align*}
These formulae are analogous to those of No.~19, and they also satisfy the relation $af + bg + ch = 0$.

22.~The formulae of Nos.~19 and 20 are connected by the principle of duality. The coordinates of a line may be regarded as point-coordinates or as plane-coordinates; and the formulae for the one case are transformed into those for the other case by the substitution of plane-coordinates for point-coordinates.

23.~If $(x, y, z, w)$ is a point on the line $(a, b, c, f, g, h)$, then the coordinates satisfy the equations
\begin{align*}
c y - b z + g w &= 0,\\
-c x + a z - f w &= 0,\\
b x - a y + h w &= 0,\\
-g x + f y - h z &= 0.
\end{align*}
These are the equations which express that the point $(x, y, z, w)$ lies on the line $(a, b, c, f, g, h)$.

\newpage
\setcounter{page}{76}

\noindent\paperref{435}\quad
\textbf{ON THE SIX COORDINATES OF A LINE.}

24.~If $(X, Y, Z, W)$ is a plane through the line $(a, b, c, f, g, h)$, then the coordinates satisfy the equations
\begin{align*}
c Y - b Z + g W &= 0,\\
-c X + a Z - f W &= 0,\\
b X - a Y + h W &= 0,\\
-g X + f Y - h Z &= 0.
\end{align*}
These are the equations which express that the plane $(X, Y, Z, W)$ passes through the line $(a, b, c, f, g, h)$.

25.~The six coordinates $(a, b, c, f, g, h)$ may be expressed in terms of four arbitrary quantities. If we write
\[
a = p q' - p' q, \quad b = r s' - r' s, \quad c = t u' - t' u,
\]
\[
f = p r' - p' r, \quad g = q s' - q' s, \quad h = t v' - t' v,
\]
where $p$, $q$, $r$, $s$, $t$, $u$, $v$, $p'$, $q'$, $r'$, $s'$, $t'$, $u'$, $v'$ are arbitrary quantities subject to the condition $af + bg + ch = 0$, then these expressions give the general solution of the equation $af + bg + ch = 0$.

26.~The general solution may also be written in the form
\begin{align*}
a &= \lambda_1 \mu_2 - \lambda_2 \mu_1, & b &= \lambda_1 \nu_2 - \lambda_2 \nu_1, & c &= \lambda_1 \rho_2 - \lambda_2 \rho_1,\\
f &= \mu_1 \nu_2 - \mu_2 \nu_1, & g &= \nu_1 \rho_2 - \nu_2 \rho_1, & h &= \rho_1 \mu_2 - \rho_2 \mu_1,
\end{align*}
where $\lambda_1$, $\mu_1$, $\nu_1$, $\rho_1$, $\lambda_2$, $\mu_2$, $\nu_2$, $\rho_2$ are arbitrary quantities.

\newpage
\setcounter{page}{77}

\noindent\paperref{435}\quad
\textbf{ON THE SIX COORDINATES OF A LINE.}

\textit{Expression for the moment of two lines.}

27.~If $(a, b, c, f, g, h)$ and $(a_1, b_1, c_1, f_1, g_1, h_1)$ are the coordinates of two lines, the moment of the two lines is defined as the product of the perpendicular distance between them into the sine of the angle between them. The expression for the moment is
\[
M = a f_1 + b g_1 + c h_1 + f a_1 + g b_1 + h c_1.
\]
This expression is symmetrical with respect to the two lines, and it vanishes when the lines intersect.

28.~The moment may also be expressed in the form
\[
M = \sqrt{(af + bg + ch)(a_1 f_1 + b_1 g_1 + c_1 h_1)} \cdot \sin\theta,
\]
where $\theta$ is the angle between the lines. But since $af + bg + ch = 0$ and $a_1 f_1 + b_1 g_1 + c_1 h_1 = 0$, this expression is indeterminate; and the proper expression is that given in No.~27.

29.~The square of the moment is given by the formula
\[
M^2 = (a f_1 + b g_1 + c h_1 + f a_1 + g b_1 + h c_1)^2.
\]
This may be expressed as a determinant; but the expression is complicated, and it is generally more convenient to use the formula of No.~27.

\newpage
\setcounter{page}{78}

\noindent\paperref{435}\quad
\textbf{ON THE SIX COORDINATES OF A LINE.}

\textit{Expression for the distance of a point from a line.}

30.~If $(x, y, z, w)$ is a point, and $(a, b, c, f, g, h)$ is a line, then the distance of the point from the line is given by the formula
\[
d = \frac{\sqrt{(c y - b z + g w)^2 + (-c x + a z - f w)^2 + (b x - a y + h w)^2}}{\sqrt{a^2 + b^2 + c^2 + f^2 + g^2 + h^2}}.
\]
This formula may be simplified when the coordinates are normalized so that $a^2 + b^2 + c^2 + f^2 + g^2 + h^2 = 1$.

31.~The distance may also be expressed in terms of the moment of the line and a line through the point parallel to the given line. If $(a_1, b_1, c_1, f_1, g_1, h_1)$ is a line through the point $(x, y, z, w)$ parallel to the line $(a, b, c, f, g, h)$, then the distance of the point from the line is equal to the moment of the two lines divided by the square root of the sum of the squares of the coordinates.

\textit{Expression for the angle between two lines.}

32.~If $(a, b, c, f, g, h)$ and $(a_1, b_1, c_1, f_1, g_1, h_1)$ are two lines, the angle $\theta$ between them is given by the formula
\[
\cos\theta = \frac{a a_1 + b b_1 + c c_1 + f f_1 + g g_1 + h h_1}{\sqrt{a^2 + b^2 + c^2 + f^2 + g^2 + h^2}\,\sqrt{a_1^2 + b_1^2 + c_1^2 + f_1^2 + g_1^2 + h_1^2}}.
\]

\newpage
\setcounter{page}{79}

\noindent\paperref{435}\quad
\textbf{ON THE SIX COORDINATES OF A LINE.}

\textit{Case of the fivefold relation.}

33.~The fivefold relation
\[
\begin{Vmatrix}
a, & b, & c, & f, & g, & h \\
a_1, & b_1, & c_1, & f_1, & g_1, & h_1 \\
a_2, & b_2, & c_2, & f_2, & g_2, & h_2 \\
a_3, & b_3, & c_3, & f_3, & g_3, & h_3 \\
a_4, & b_4, & c_4, & f_4, & g_4, & h_4
\end{Vmatrix}
= 0
\]
expresses that the line $(a, b, c, f, g, h)$ belongs to the linear complex determined by the four lines $(a_1, b_1, \ldots)$, $(a_2, b_2, \ldots)$, $(a_3, b_3, \ldots)$, $(a_4, b_4, \ldots)$. This relation may be regarded as the condition that the five lines have a common transversal; or that they belong to the same linear congruence.

34.~The fivefold relation may be written in the form of a determinant of the fifth order; but it is generally more convenient to use the form given in No.~33. The relation is equivalent to five linear equations in the coordinates $(a, b, c, f, g, h)$, which express that the line belongs to the linear complex determined by the four given lines.

\textit{Case of the fourfold relation.}

35.~The fourfold relation
\[
\begin{Vmatrix}
a, & b, & c, & f, & g, & h \\
a_1, & b_1, & c_1, & f_1, & g_1, & h_1 \\
a_2, & b_2, & c_2, & f_2, & g_2, & h_2 \\
a_3, & b_3, & c_3, & f_3, & g_3, & h_3
\end{Vmatrix}
= 0
\]
expresses that the line $(a, b, c, f, g, h)$ belongs to the regulus determined by the three lines $(a_1, b_1, \ldots)$, $(a_2, b_2, \ldots)$, $(a_3, b_3, \ldots)$. This relation is equivalent to four linear equations in the coordinates $(a, b, c, f, g, h)$, which express that the line meets each of the three given lines.

\newpage
\setcounter{page}{80}

\noindent\paperref{435}\quad
\textbf{ON THE SIX COORDINATES OF A LINE.}

\smallskip
\textbf{Article, Nos.~36 to 58. Further Developments.}

\smallskip
\textit{Expression for the coordinates of the line of intersection of two planes.}

36.~If $(X, Y, Z, W)$ and $(X', Y', Z', W')$ are the coordinates of two planes, the six coordinates of their line of intersection are given by the formulae
\begin{align*}
a &= X W' - X' W, & b &= Y W' - Y' W, & c &= Z W' - Z' W,\\
f &= Y Z' - Y' Z, & g &= Z X' - Z' X, & h &= X Y' - X' Y;
\end{align*}
and these satisfy the relation $af + bg + ch = 0$.

37.~The coordinates of the line joining two points $(x, y, z, w)$ and $(x', y', z', w')$ are given by the formulae
\begin{align*}
a &= y z' - y' z, & b &= z x' - z' x, & c &= x y' - x' y,\\
f &= x w' - x' w, & g &= y w' - y' w, & h &= z w' - z' w.
\end{align*}

38.~The condition that the line $(a, b, c, f, g, h)$ may meet the line $(a_1, b_1, c_1, f_1, g_1, h_1)$ is
\[
a f_1 + b g_1 + c h_1 + f a_1 + g b_1 + h c_1 = 0.
\]
This condition may be written in the form of a determinant; but the form given above is generally more convenient.

\newpage
\setcounter{page}{81}

\noindent\paperref{435}\quad
\textbf{ON THE SIX COORDINATES OF A LINE.}

\textit{Expression for the coordinates of a line in terms of six arbitrary quantities.}

39.~The six coordinates $(a, b, c, f, g, h)$ of a line are connected by the single equation $af + bg + ch = 0$; they may therefore be expressed in terms of five independent quantities. If we introduce a sixth quantity $k$, and write
\begin{align*}
a &= k \alpha, & b &= k \beta, & c &= k \gamma,\\
f &= k \phi, & g &= k \psi, & h &= k \chi,
\end{align*}
where $\alpha$, $\beta$, $\gamma$, $\phi$, $\psi$, $\chi$ are quantities satisfying the equation $\alpha\phi + \beta\psi + \gamma\chi = 0$, then the ratios $a : b : c : f : g : h$ are determined by the five quantities $\alpha : \beta : \gamma : \phi : \psi : \chi$.

40.~The absolute values of the coordinates may be fixed by imposing an additional condition; for example, we may assume that $a^2 + b^2 + c^2 + f^2 + g^2 + h^2 = 1$, or that one of the coordinates is equal to unity. But it is generally more convenient to leave the absolute values indeterminate, and to consider only the ratios of the coordinates.

41.~The coordinates of a line may be normalized so that $a^2 + b^2 + c^2 + f^2 + g^2 + h^2 = 1$; and when this is done, the distance of a point from the line, and the angle between two lines, may be expressed in simpler forms. But the normalization is not essential, and it is often more convenient to use the unnormalized coordinates.

\newpage
\setcounter{page}{82}

\noindent\paperref{435}\quad
\textbf{ON THE SIX COORDINATES OF A LINE.}

\textit{Transformation of coordinates.}

42.~If we change the system of coordinate planes, the six coordinates of a line undergo a linear transformation. Let the new coordinate planes be $x_0 = 0$, $y_0 = 0$, $z_0 = 0$, $w_0 = 0$, where
\begin{align*}
x_0 &= a_1 x + b_1 y + c_1 z + d_1 w,\\
y_0 &= a_2 x + b_2 y + c_2 z + d_2 w,\\
z_0 &= a_3 x + b_3 y + c_3 z + d_3 w,\\
w_0 &= a_4 x + b_4 y + c_4 z + d_4 w.
\end{align*}
Then if $(a, b, c, f, g, h)$ are the coordinates of a line referred to the old system, and $(a_0, b_0, c_0, f_0, g_0, h_0)$ the coordinates of the same line referred to the new system, we have
\begin{align*}
a_0 &= (a_2 b_3 - a_3 b_2) a + (a_3 b_1 - a_1 b_3) b + (a_1 b_2 - a_2 b_1) c\\
&\quad + (a_2 c_3 - a_3 c_2) f + (a_3 c_1 - a_1 c_3) g + (a_1 c_2 - a_2 c_1) h,\\
&\ \,\vdots
\end{align*}
with similar formulae for $b_0$, $c_0$, $f_0$, $g_0$, $h_0$.

43.~These formulae of transformation are linear in the coordinates $(a, b, c, f, g, h)$, and they preserve the relation $af + bg + ch = 0$. They may be written in the matrix form
\[
\begin{pmatrix} a_0 \\ b_0 \\ c_0 \\ f_0 \\ g_0 \\ h_0 \end{pmatrix}
= M
\begin{pmatrix} a \\ b \\ c \\ f \\ g \\ h \end{pmatrix},
\]
where $M$ is a $6 \times 6$ matrix whose elements are the minors of the matrix of transformation of point-coordinates.

\newpage
\setcounter{page}{83}

\noindent\paperref{435}\quad
\textbf{ON THE SIX COORDINATES OF A LINE.}

44.~The matrix $M$ satisfies the condition that it preserves the bilinear form $af + bg + ch$; that is, if $(a, b, c, f, g, h)$ satisfy $af + bg + ch = 0$, then $(a_0, b_0, c_0, f_0, g_0, h_0)$ also satisfy $a_0 f_0 + b_0 g_0 + c_0 h_0 = 0$. This condition is expressed by the equation
\[
M^T J M = J,
\]
where $J$ is the matrix
\[
J = \begin{pmatrix}
0 & 0 & 0 & 1 & 0 & 0 \\
0 & 0 & 0 & 0 & 1 & 0 \\
0 & 0 & 0 & 0 & 0 & 1 \\
1 & 0 & 0 & 0 & 0 & 0 \\
0 & 1 & 0 & 0 & 0 & 0 \\
0 & 0 & 1 & 0 & 0 & 0
\end{pmatrix}.
\]

45.~The transformation formulae may also be expressed in terms of the coordinates of two planes. If $(X, Y, Z, W)$ and $(X', Y', Z', W')$ are the coordinates of two planes, and $(X_0, Y_0, Z_0, W_0)$, $(X'_0, Y'_0, Z'_0, W'_0)$ are their coordinates referred to the new system, then the coordinates of the line of intersection are transformed by the same formulae as above.

46.~The absolute magnitudes of the coordinates may be fixed by imposing the condition that the determinant of the matrix of transformation shall be equal to unity. When this is done, the coordinates are said to be \textit{normalized}; and the formulae of transformation are then completely determined.

\newpage
\setcounter{page}{84}

\noindent\paperref{435}\quad
\textbf{ON THE SIX COORDINATES OF A LINE.}

\textit{Condition that four lines may belong to the same regulus.}

47.~If $(a_1, b_1, c_1, f_1, g_1, h_1)$, $(a_2, b_2, c_2, f_2, g_2, h_2)$, $(a_3, b_3, c_3, f_3, g_3, h_3)$, $(a_4, b_4, c_4, f_4, g_4, h_4)$ are four lines, the condition that they may belong to the same regulus is that the determinant
\[
\begin{vmatrix}
a_1 & b_1 & c_1 & f_1 & g_1 & h_1 \\
a_2 & b_2 & c_2 & f_2 & g_2 & h_2 \\
a_3 & b_3 & c_3 & f_3 & g_3 & h_3 \\
a_4 & b_4 & c_4 & f_4 & g_4 & h_4
\end{vmatrix}
\]
shall be of rank not exceeding 3; that is, all the determinants of the fourth order formed from it shall vanish.

48.~This condition is equivalent to the three conditions that the four lines shall have a common transversal; or that they shall belong to the same linear congruence. It may also be expressed as the condition that the six coordinates of any one of the lines can be expressed as linear functions of the coordinates of the other three.

49.~The condition that four lines may belong to the same regulus may be written in the form
\[
\begin{vmatrix}
a_1 f_2 + b_1 g_2 + c_1 h_2 + f_1 a_2 + g_1 b_2 + h_1 c_2 & \cdots \\ \vdots & \ddots
\end{vmatrix}
= 0,
\]
where the determinant is of the fourth order, and its elements are the moments of pairs of lines.

\newpage
\setcounter{page}{85}

\noindent\paperref{435}\quad
\textbf{ON THE SIX COORDINATES OF A LINE.}

\textit{Condition that five lines may have a common transversal.}

50.~If $(a_1, b_1, c_1, f_1, g_1, h_1)$, $(a_2, b_2, c_2, f_2, g_2, h_2)$, $(a_3, b_3, c_3, f_3, g_3, h_3)$, $(a_4, b_4, c_4, f_4, g_4, h_4)$, $(a_5, b_5, c_5, f_5, g_5, h_5)$ are five lines, the condition that they may have a common transversal is that the determinant
\[
\begin{vmatrix}
a_1 & b_1 & c_1 & f_1 & g_1 & h_1 \\
a_2 & b_2 & c_2 & f_2 & g_2 & h_2 \\
a_3 & b_3 & c_3 & f_3 & g_3 & h_3 \\
a_4 & b_4 & c_4 & f_4 & g_4 & h_4 \\
a_5 & b_5 & c_5 & f_5 & g_5 & h_5
\end{vmatrix}
= 0,
\]
where the determinant is of the fifth order. This condition is equivalent to the condition that the five lines belong to the same linear complex; or that they have $\infty^1$ common transversals.

51.~The condition may also be expressed as the vanishing of the determinant of the matrix whose rows are the coordinates of the five lines. This determinant is of the fifth order, and its vanishing expresses that the five lines are linearly dependent.

\textit{Condition that six lines may have two common transversals.}

52.~If $(a_1, b_1, c_1, f_1, g_1, h_1)$, $\ldots$, $(a_6, b_6, c_6, f_6, g_6, h_6)$ are six lines, the condition that they may have two common transversals is that the determinant
\[
\begin{vmatrix}
a_1 & b_1 & c_1 & f_1 & g_1 & h_1 \\
\vdots & \vdots & \vdots & \vdots & \vdots & \vdots \\
a_6 & b_6 & c_6 & f_6 & g_6 & h_6
\end{vmatrix}
= 0,
\]
where the determinant is of the sixth order. This condition expresses that the six lines belong to the same linear congruence; or that they have $\infty^2$ common transversals.

\newpage
\setcounter{page}{86}

\noindent\paperref{435}\quad
\textbf{ON THE SIX COORDINATES OF A LINE.}

53.~The condition that six lines may have two common transversals may also be expressed as the condition that the six lines are linearly dependent; that is, that there exist quantities $\lambda_1$, $\lambda_2$, $\ldots$, $\lambda_6$, not all zero, such that
\[
\lambda_1 (a_1, b_1, c_1, f_1, g_1, h_1) + \cdots + \lambda_6 (a_6, b_6, c_6, f_6, g_6, h_6) = (0, 0, 0, 0, 0, 0).
\]
This condition is equivalent to the vanishing of the determinant of the sixth order formed from the coordinates of the six lines.

\textit{Expression for the coordinates of the common transversal of four lines.}

54.~If four lines are given, the coordinates of their common transversal may be found by solving the four linear equations which express that the transversal meets each of the given lines. Let $(a, b, c, f, g, h)$ be the coordinates of the transversal; then the conditions that it meets the four given lines are
\begin{align*}
a f_1 + b g_1 + c h_1 + f a_1 + g b_1 + h c_1 &= 0,\\
a f_2 + b g_2 + c h_2 + f a_2 + g b_2 + h c_2 &= 0,\\
a f_3 + b g_3 + c h_3 + f a_3 + g b_3 + h c_3 &= 0,\\
a f_4 + b g_4 + c h_4 + f a_4 + g b_4 + h c_4 &= 0.
\end{align*}
These four equations, together with $af + bg + ch = 0$, determine the ratios of the six coordinates.

55.~The coordinates of the transversal may be expressed as determinants of the fourth order formed from the coordinates of the four given lines. If we write
\[
A = \begin{vmatrix} f_1 & g_1 & h_1 & a_1 \\ \vdots & & & \\ f_4 & g_4 & h_4 & a_4 \end{vmatrix},
\quad \text{etc.},
\]
then the coordinates of the transversal are proportional to $(A, B, C, F, G, H)$.

\newpage
\setcounter{page}{87}

\noindent\paperref{435}\quad
\textbf{ON THE SIX COORDINATES OF A LINE.}

56.~The coordinates of the transversal may also be found by the following construction. Let the four given lines be $L_1$, $L_2$, $L_3$, $L_4$; through $L_1$ and $L_2$ draw two arbitrary planes, and let their line of intersection be $M_{12}$; through $L_3$ and $L_4$ draw two arbitrary planes, and let their line of intersection be $M_{34}$. Then the common transversal of $L_1$, $L_2$, $L_3$, $L_4$ is the line of intersection of the two planes which contain $M_{12}$ and $M_{34}$ respectively, and which pass through the line of intersection of the planes of $L_1$, $L_2$ and $L_3$, $L_4$.

57.~This construction gives a geometrical method of finding the common transversal of four lines; and it may be extended to the case of five or more lines. The construction is based on the principle that the common transversal of four lines is the line of intersection of two planes, each of which contains two of the given lines.

\textit{Expression for the coordinates of the line of shortest distance between two lines.}

58.~If $(a, b, c, f, g, h)$ and $(a_1, b_1, c_1, f_1, g_1, h_1)$ are two lines, the coordinates of their common perpendicular (or line of shortest distance) may be found as follows. Let $(\lambda, \mu, \nu, \rho, \sigma, \tau)$ be the coordinates of the common perpendicular; then the conditions that it meets each of the two lines are
\begin{align*}
\lambda f + \mu g + \nu h + \rho a + \sigma b + \tau c &= 0,\\
\lambda f_1 + \mu g_1 + \nu h_1 + \rho a_1 + \sigma b_1 + \tau c_1 &= 0.
\end{align*}
And the condition that it is perpendicular to each of the two lines is that its direction-cosines are proportional to the differences of the direction-cosines of the two lines.

\newpage
\setcounter{page}{88}

\noindent\paperref{435}\quad
\textbf{ON THE SIX COORDINATES OF A LINE.}

\smallskip
\textbf{Article, Nos.~59 to 75. Statical and Kinematical Applications.}

\smallskip
The coordinates $(a, b, c, f, g, h)$, as last defined, are peculiarly convenient in kinematical and mechanical questions, as will appear from the following investigations.

\textit{Rotations.}

59.~Using the term rotation to denote an infinitesimal rotation, I say first that a rotation $\lambda$ round the line $(a, b, c, f, g, h)$ produces in the point $(x, y, z)$ rigidly connected with this line the displacements
\begin{align*}
\delta x &= \lambda(\phantom{hx} - hy + gz - a),\\
\delta y &= \lambda(\phantom{hx} hx \phantom{+} - fz - b),\\
\delta z &= \lambda(-gx + fy \phantom{+ gz} - c).
\end{align*}

\textit{Forces.}

60.~A force $R$ acting along the line $(a, b, c, f, g, h)$ has for its components along the axes
\begin{align*}
X &= R a, & Y &= R b, & Z &= R c,
\end{align*}
and its moments about the axes are
\begin{align*}
L &= R f, & M &= R g, & N &= R h.
\end{align*}
These six quantities satisfy the relation $a f + b g + c h = 0$, which is the condition that the force may be a single force acting along a definite line.

\newpage
\setcounter{page}{89}

\noindent\paperref{435}\quad
\textbf{ON THE SIX COORDINATES OF A LINE.}

\textit{Couples.}

61.~A couple $G$ in the plane $(l, m, n, p)$ has for its components
\begin{align*}
L &= G l, & M &= G m, & N &= G n,
\end{align*}
and its forces along the axes are
\begin{align*}
X &= G p, & Y &= G q, & Z &= G r,
\end{align*}
where $p$, $q$, $r$ are determined by the condition that the system is equivalent to a couple.

62.~The six components of a couple satisfy the relation $a f + b g + c h = 0$, which is the condition that the system may be reducible to a single force or a couple. In fact, the general system of forces acting on a rigid body may be reduced to a single force acting along a certain line, together with a couple in a plane perpendicular to that line.

\textit{Wrenches.}

63.~A \textit{wrench} is a force $R$ acting along a line $(a, b, c, f, g, h)$, together with a couple $G$ in a plane perpendicular to the line. The six components of the wrench are
\begin{align*}
X &= R a, & Y &= R b, & Z &= R c,\\
L &= R f + G a, & M &= R g + G b, & N &= R h + G c.
\end{align*}
These six quantities are connected by the relation
\[
a L + b M + c N - f X - g Y - h Z = 0,
\]
which is the condition that the system may be a wrench.

\newpage
\setcounter{page}{90}

\noindent\paperref{435}\quad
\textbf{ON THE SIX COORDINATES OF A LINE.}

64.~The pitch of the wrench is defined as the ratio $p = G/R$ of the couple to the force. The six components may then be written
\begin{align*}
X &= R a, & Y &= R b, & Z &= R c,\\
L &= R(f + pa), & M &= R(g + pb), & N &= R(h + pc).
\end{align*}
And the relation between them becomes
\[
a L + b M + c N - f X - g Y - h Z = p(a X + b Y + c Z) = 0,
\]
since $a X + b Y + c Z = R(a^2 + b^2 + c^2)$.

65.~The axis of the wrench is the line $(a, b, c, f, g, h)$; and the pitch is the ratio of the couple to the force. The wrench is completely determined by its axis and its pitch; and conversely, any line and any pitch determine a wrench.

66.~If the pitch is zero, the wrench reduces to a single force; if the pitch is infinite, the wrench reduces to a couple. The general wrench is therefore intermediate between a force and a couple.

\textit{Screws.}

67.~A \textit{screw} is a line associated with a pitch. The six coordinates of a screw are the six coordinates of the line, together with the pitch. The screw may be regarded as a wrench of unit force; and its six components are
\begin{align*}
X &= a, & Y &= b, & Z &= c,\\
L &= f + pa, & M &= g + pb, & N &= h + pc.
\end{align*}

\newpage
\setcounter{page}{91}

\noindent\paperref{435}\quad
\textbf{ON THE SIX COORDINATES OF A LINE.}

68.~The theory of screws was developed by Sir Robert Ball in his ``Theory of Screws'' (1876). A screw is determined by its axis $(a, b, c, f, g, h)$ and its pitch $p$; and it may be regarded as a geometrical entity consisting of a line and a scalar quantity associated with it.

69.~Two screws are said to be \textit{reciprocal} when the virtual coefficient of the one with respect to the other vanishes. If $(a, b, c, f, g, h, p)$ and $(a_1, b_1, c_1, f_1, g_1, h_1, p_1)$ are two screws, the condition of reciprocity is
\[
a f_1 + b g_1 + c h_1 + f a_1 + g b_1 + h c_1 + p(a a_1 + b b_1 + c c_1) + p_1(a_1 a + b_1 b + c_1 c) = 0.
\]

70.~A system of screws is said to be a \textit{co-reciprocal system} when every pair of screws in the system are reciprocal. Such a system contains at most six screws; and when it contains exactly six, the screws are said to form a \textit{screw-six}.

\textit{Applications to Mechanics.}

71.~The six coordinates of a line are of great importance in the theory of the equilibrium and motion of rigid bodies. The condition of equilibrium of a rigid body under the action of any system of forces is that the sum of the components and the sum of the moments about any line shall vanish.

\newpage
\setcounter{page}{92}

\noindent\paperref{435}\quad
\textbf{ON THE SIX COORDINATES OF A LINE.}

72.~If $(X_i, Y_i, Z_i, L_i, M_i, N_i)$ are the components of the $i$th force, then the conditions of equilibrium are
\begin{align*}
\sum X_i &= 0, & \sum Y_i &= 0, & \sum Z_i &= 0,\\
\sum L_i &= 0, & \sum M_i &= 0, & \sum N_i &= 0.
\end{align*}
These six equations express that the resultant force and the resultant couple vanish.

73.~The six coordinates of a line are also useful in the theory of virtual work. If a rigid body undergoes a small displacement, the work done by the forces acting on it is equal to the sum of the products of the components of the forces into the corresponding displacements.

74.~If $(\delta x, \delta y, \delta z)$ is the displacement of a point, and $(\delta\theta_x, \delta\theta_y, \delta\theta_z)$ is the rotation of the body, then the virtual work is
\[
\sum (X_i \delta x + Y_i \delta y + Z_i \delta z + L_i \delta\theta_x + M_i \delta\theta_y + N_i \delta\theta_z) = 0.
\]
This equation expresses the principle of virtual work for a rigid body.

75.~The six coordinates of a line are also applicable to the theory of attractions and potentials. If a body is attracted by a system of forces, the potential of the system may be expressed in terms of the six coordinates of the lines of action of the forces.

\newpage
\setcounter{page}{93}

\noindent\paperref{435}\quad
\textbf{ON THE SIX COORDINATES OF A LINE.}

\smallskip
\textbf{Article, Nos.~76 to 99. Further Analytical Developments.}

\smallskip
\textit{Expression for the condition that six lines may be generators of the same quadric surface.}

76.~If six lines are generators of the same quadric surface, their coordinates satisfy a certain condition. Let the equation of the quadric surface be
\[
Ax^2 + By^2 + Cz^2 + Dw^2 + 2Fyz + 2Gzx + 2Hxy + 2Lxw + 2Myw + 2Nzw = 0.
\]
Then the condition that the line $(a, b, c, f, g, h)$ may be a generator of the quadric is
\[
A a^2 + B b^2 + C c^2 + D f^2 + E g^2 + F h^2 + 2F bc + 2G ca + 2H ab + 2L af + 2M bg + 2N ch = 0.
\]
This is a quadratic equation in the coordinates of the line; and the condition that six lines may be generators of the same quadric is that their coordinates satisfy this equation.

77.~The condition may be written in the form
\[
\sum_{i,j} Q_{ij} a_i a_j = 0,
\]
where $Q_{ij}$ are the coefficients of the quadric, and $a_i$ are the coordinates of the line. This is a quadratic condition; and the six lines satisfy it when they are generators of the same quadric.

\newpage
\setcounter{page}{94}

\noindent\paperref{435}\quad
\textbf{ON THE SIX COORDINATES OF A LINE.}

\textit{Expression for the condition that a line may touch a given quadric surface.}

78.~If the quadric surface is given by the equation
\[
Ax^2 + By^2 + Cz^2 + Dw^2 + 2Fyz + 2Gzx + 2Hxy + 2Lxw + 2Myw + 2Nzw = 0,
\]
then the condition that the line $(a, b, c, f, g, h)$ may touch the quadric is that the discriminant of the quadratic form vanish. This condition is
\[
\begin{vmatrix}
A & H & G & L \\
H & B & F & M \\
G & F & C & N \\
L & M & N & D
\end{vmatrix}
= 0.
\]
This is a quartic condition in the coordinates of the line.

79.~The condition may also be expressed as the condition that the line is a tangent to the quadric; that is, that it meets the quadric in two coincident points. This condition is equivalent to the condition that the line is a generator of the quadric, or that it lies altogether on the quadric.

\textit{Expression for the condition that a line may belong to a given linear complex.}

80.~The condition that the line $(a, b, c, f, g, h)$ may belong to the linear complex
\[
AX + BY + CZ + DW + EU + FV = 0
\]
is
\[
Aa + Bb + Cc + Df + Eg + Fh = 0.
\]
This is a linear condition in the coordinates of the line; and the lines belonging to a linear complex form a threefold system.

\newpage
\setcounter{page}{95}

\noindent\paperref{435}\quad
\textbf{ON THE SIX COORDINATES OF A LINE.}

81.~The linear complex may be represented as the assemblage of lines which satisfy the linear equation
\[
Aa + Bb + Cc + Df + Eg + Fh = 0.
\]
The coefficients $(A, B, C, D, E, F)$ are the coordinates of the complex; and the equation expresses that the line $(a, b, c, f, g, h)$ belongs to the complex.

82.~The condition that a line may belong to a given linear congruence is that its coordinates satisfy two linear equations
\begin{align*}
A_1 a + B_1 b + C_1 c + D_1 f + E_1 g + F_1 h &= 0,\\
A_2 a + B_2 b + C_2 c + D_2 f + E_2 g + F_2 h &= 0.
\end{align*}
The lines belonging to a linear congruence form a twofold system; they are the common lines of two linear complexes.

83.~The condition that a line may belong to a given regulus is that its coordinates satisfy three linear equations. The lines belonging to a regulus form a onefold system; they are the generators of a quadric surface.

\textit{Expression for the condition that four lines may have a common transversal.}

84.~If $(a_1, b_1, c_1, f_1, g_1, h_1)$, $(a_2, b_2, c_2, f_2, g_2, h_2)$, $(a_3, b_3, c_3, f_3, g_3, h_3)$, $(a_4, b_4, c_4, f_4, g_4, h_4)$ are four lines, the condition that they may have a common transversal is that the determinant
\[
\begin{vmatrix}
a_1 & b_1 & c_1 & f_1 & g_1 & h_1 \\
a_2 & b_2 & c_2 & f_2 & g_2 & h_2 \\
a_3 & b_3 & c_3 & f_3 & g_3 & h_3 \\
a_4 & b_4 & c_4 & f_4 & g_4 & h_4
\end{vmatrix}
= 0.
\]

\newpage
\setcounter{page}{96}

\noindent\paperref{435}\quad
\textbf{ON THE SIX COORDINATES OF A LINE.}

This condition is equivalent to the condition that the four lines belong to the same linear congruence; or that the six coordinates of their common transversal satisfy the four linear equations which express that the transversal meets each of the given lines.

85.~The condition may also be expressed as the vanishing of the determinant of the matrix whose rows are the coordinates of the four lines. This determinant is of the fourth order in the elements of each line; and its vanishing expresses that the four lines are linearly dependent.

86.~The condition that four lines may have a common transversal is equivalent to the condition that the four lines belong to the same regulus; for if four lines have a common transversal, they belong to the linear congruence determined by the transversal and any line meeting it; and conversely, if four lines belong to the same regulus, they have a common transversal.

\textit{Expression for the condition that five lines may belong to the same linear complex.}

87.~If $(a_1, b_1, c_1, f_1, g_1, h_1)$, $\ldots$, $(a_5, b_5, c_5, f_5, g_5, h_5)$ are five lines, the condition that they may belong to the same linear complex is that the determinant
\[
\begin{vmatrix}
a_1 & b_1 & c_1 & f_1 & g_1 & h_1 \\
\vdots & \vdots & \vdots & \vdots & \vdots & \vdots \\
a_5 & b_5 & c_5 & f_5 & g_5 & h_5
\end{vmatrix}
= 0.
\]

\newpage
\setcounter{page}{97}

\noindent\paperref{435}\quad
\textbf{ON THE SIX COORDINATES OF A LINE.}

This condition is equivalent to the condition that the five lines are linearly dependent; that is, that there exist quantities $\lambda_1$, $\ldots$, $\lambda_5$, not all zero, such that
\[
\lambda_1 (a_1, b_1, c_1, f_1, g_1, h_1) + \cdots + \lambda_5 (a_5, b_5, c_5, f_5, g_5, h_5) = (0, 0, 0, 0, 0, 0).
\]

88.~The condition may also be expressed as the vanishing of the determinant of the fifth order formed from the coordinates of the five lines. This determinant is a skew-symmetric determinant of the fifth order; and its vanishing expresses that the five lines belong to the same linear complex.

89.~The condition that five lines may belong to the same linear complex is equivalent to the condition that they have $\infty^1$ common transversals; for if five lines belong to a linear complex, any line of the complex meets them all; and conversely, if five lines have $\infty^1$ common transversals, these transversals form a linear complex, and the five lines belong to it.

\textit{Expression for the condition that six lines may belong to the same linear congruence.}

90.~If $(a_1, b_1, c_1, f_1, g_1, h_1)$, $\ldots$, $(a_6, b_6, c_6, f_6, g_6, h_6)$ are six lines, the condition that they may belong to the same linear congruence is that the determinant
\[
\begin{vmatrix}
a_1 & b_1 & c_1 & f_1 & g_1 & h_1 \\
\vdots & \vdots & \vdots & \vdots & \vdots & \vdots \\
a_6 & b_6 & c_6 & f_6 & g_6 & h_6
\end{vmatrix}
= 0.
\]

\newpage
\setcounter{page}{98}

\noindent\paperref{435}\quad
\textbf{ON THE SIX COORDINATES OF A LINE.}

This condition is equivalent to the condition that the six lines are linearly dependent; that is, that there exist quantities $\lambda_1$, $\ldots$, $\lambda_6$, not all zero, such that
\[
\lambda_1 (a_1, b_1, c_1, f_1, g_1, h_1) + \cdots + \lambda_6 (a_6, b_6, c_6, f_6, g_6, h_6) = (0, 0, 0, 0, 0, 0).
\]

91.~The condition may also be expressed as the vanishing of the determinant of the sixth order formed from the coordinates of the six lines. This determinant is a skew-symmetric determinant of the sixth order; and its vanishing expresses that the six lines belong to the same linear congruence.

92.~The condition that six lines may belong to the same linear congruence is equivalent to the condition that they have $\infty^2$ common transversals; for if six lines belong to a linear congruence, any line of the congruence meets them all; and conversely, if six lines have $\infty^2$ common transversals, these transversals form a linear congruence, and the six lines belong to it.

\textit{Expression for the condition that seven lines may have a common transversal.}

93.~If seven lines are given, the condition that they may have a common transversal is that the coordinates of the lines satisfy a certain system of equations. This condition is equivalent to the condition that the seven lines belong to the same linear complex; or that they have $\infty^1$ common transversals.

\newpage
\setcounter{page}{99}

\noindent\paperref{435}\quad
\textbf{ON THE SIX COORDINATES OF A LINE.}

94.~The condition that seven lines may have a common transversal may be expressed as the vanishing of certain determinants formed from the coordinates of the lines. If $(a_1, b_1, c_1, f_1, g_1, h_1)$, $\ldots$, $(a_7, b_7, c_7, f_7, g_7, h_7)$ are the seven lines, then the condition is that all the determinants of the sixth order formed from the matrix
\[
\begin{pmatrix}
a_1 & b_1 & c_1 & f_1 & g_1 & h_1 \\
\vdots & \vdots & \vdots & \vdots & \vdots & \vdots \\
a_7 & b_7 & c_7 & f_7 & g_7 & h_7
\end{pmatrix}
\]
shall vanish. This condition is equivalent to seven independent equations; and when it is satisfied, the seven lines have a common transversal.

95.~The condition may also be expressed as the condition that the seven lines are linearly dependent; that is, that there exist quantities $\lambda_1$, $\ldots$, $\lambda_7$, not all zero, such that
\[
\lambda_1 (a_1, b_1, c_1, f_1, g_1, h_1) + \cdots + \lambda_7 (a_7, b_7, c_7, f_7, g_7, h_7) = (0, 0, 0, 0, 0, 0).
\]

\bigskip
\noindent\rule{\textwidth}{0.4pt}

% ============================================================
% PAPER 436: ON A CERTAIN SEXTIC TORSE
% ============================================================
\newpage
\setcounter{page}{99}

\noindent\paperref{436}\quad
\begin{center}
\textbf{436.}\\[6pt]
\textbf{ON A CERTAIN SEXTIC TORSE.}\\[4pt]
\small
\textit{[From the \textit{Transactions of the Cambridge Philosophical Society}, vol.~\textsc{xi}. Part~\textsc{iii}. (1871),\\
pp.~507---523. Read Nov.~8, 1869.]}
\end{center}

\normalsize
\setlength{\parindent}{1em}
The torse (developable surface) intended to be considered is that which has for its edge of regression an excubo-quartic curve, or say a \textit{unicursal} quartic curve. I call to mind that (excluding the plane quartic) a quartic curve is either a quadriquadric, viz. it is the complete intersection of two quadric surfaces; or else it is an excubo-quartic, viz. there is through the curve only one quadric surface, and the curve is the partial intersection of this quadric surface with a cubic surface through two generating lines (of the same kind) of the quadric surface. Returning to the quadriquadric curve, this may be general, nodal, or cuspidal; viz. if the two quadric surfaces have an ordinary contact, the curve of intersection is a nodal quadriquadric; if they have a stationary contact, the curve is a cuspidal quadriquadric.

The unicursal quartic is a curve such that the coordinates $(x, y, z, w)$ of any point thereof are proportional to rational and integral quartic functions $(*\!\!\backslash\!\!\backslash\, \theta, 1)^4$ of a variable parameter $\theta$; and the general unicursal quartic is in fact the excubo-quartic; but included as particular cases of the unicursal curve (although not as cases of the excubo-quartic as above defined) we have the nodal quadriquadric and the cuspidal quadriquadric. The torse having for its edge of regression a unicursal curve is a sextic torse; and this is in fact the order of the torse derived from the excubo-quartic, and from the nodal quadriquadric; but for the cuspidal quadriquadric, there is a depression of one, and the torse becomes a quintic torse. The equations have been obtained of (1) the sextic torse derived from the nodal quadriquadric, (2) the quintic torse derived from the cuspidal quadriquadric, (3) the sextic torse derived from a certain special excubo-quartic; but the equation of the torse derived from the general unicursal quartic has not yet been found. To show at the outset what the analytical problem is, I

\newpage
\setcounter{page}{100}

\noindent\paperref{436}\quad
\textbf{ON A CERTAIN SEXTIC TORSE.}

remark that the unicursal quartic may be represented by the equations
\begin{align*}
x : y : z : w &= (\theta - \alpha)(\theta - \beta)(\theta - \gamma)(\theta - \delta) \\
& : (\theta - \alpha')(\theta - \beta')(\theta - \gamma')(\theta - \delta') \\
& : (\theta - \alpha'')(\theta - \beta'')(\theta - \gamma'')(\theta - \delta'') \\
& : (\theta - \alpha''')(\theta - \beta''')(\theta - \gamma''')(\theta - \delta'''),
\end{align*}
where $\alpha$, $\beta$, $\gamma$, $\delta$, $\alpha'$, $\beta'$, $\gamma'$, $\delta'$, $\alpha''$, $\beta''$, $\gamma''$, $\delta''$, $\alpha'''$, $\beta'''$, $\gamma'''$, $\delta'''$ are arbitrary constants; or, what is the same thing, the curve is the partial intersection of the quadric surface
\[
\frac{x}{(\theta - \alpha)(\theta - \beta)(\theta - \gamma)(\theta - \delta)}
= \frac{y}{(\theta - \alpha')(\theta - \beta')(\theta - \gamma')(\theta - \delta')}
\]
with the cubic surface which is the locus of the point of intersection of the three planes
\begin{align*}
\frac{x}{(\theta - \alpha)(\theta - \beta)(\theta - \gamma)(\theta - \delta)} &=
\frac{z}{(\theta - \alpha'')(\theta - \beta'')(\theta - \gamma'')(\theta - \delta'')},\\
\frac{y}{(\theta - \alpha')(\theta - \beta')(\theta - \gamma')(\theta - \delta')} &=
\frac{w}{(\theta - \alpha''')(\theta - \beta''')(\theta - \gamma''')(\theta - \delta''')},\\
\frac{x}{(\theta - \alpha)(\theta - \beta)(\theta - \gamma)(\theta - \delta)} &=
\frac{w}{(\theta - \alpha''')(\theta - \beta''')(\theta - \gamma''')(\theta - \delta''')}.
\end{align*}
The torse is the envelope of the osculating plane of the curve; or, what is the same thing, it is the locus of the tangent lines to the curve. The equation of the torse is to be obtained by eliminating $\theta$ between the equations of the tangent line and the condition that the line is tangent to the curve.

\end{document}
